\documentclass[11pt,a4paper]{article}
\usepackage[french]{babel}
\usepackage[T1]{fontenc}
\usepackage[utf8]{inputenc}
\usepackage{geometry}
\geometry{margin=2.5cm}
\usepackage{hyperref}
\usepackage{xcolor}

\definecolor{larcblue}{HTML}{1565C0}

\title{\color{larcblue}Supervision des présences et événements — Documentation technique}
\author{Les logiciels Larc — Arc-en-Ciel}
\date{\today}

\begin{document}
\maketitle
\tableofcontents
\newpage

\section{Architecture}

LarcSuperviseur est une application \textbf{PySide6 (Qt6)} desktop,
développée en Python~3. Elle utilise les widgets Material Design~3 via
\texttt{phibuilder} et se connecte à PostgreSQL via \texttt{psycopg2}.

\section{Fonctionnalités principales}

\begin{itemize}
  \item Suivi des absences et retards en temps réel
  \item Génération d'événements (absence journée, retard, événements)
  \item Statistiques par groupe, classe et élève
  \item Tableau de bord avec KPIs et graphiques
  \item Gestion des photos élèves
  \item Éditeur d'emploi du temps
\end{itemize}

\section{Rôles utilisateurs}

\begin{itemize}
  \item \textbf{Superviseur (écriture)}
  \item \textbf{Coordinateur (valider)}
  \item \textbf{Administrateur (complet)}
\end{itemize}

\section{Base de données}

Tables principales~:

\begin{itemize}
  \item \texttt{student\_event}
  \item \texttt{larcauth\_type\_event}
  \item \texttt{larcauth\_lieu}
  \item \texttt{larcauth\_academicyear}
\end{itemize}

\section{Modules principaux}

\begin{center}
\begin{tabular}{|l|l|}
\hline
\textbf{Fichier} & \textbf{Rôle} \\ \hline
    \texttt{main\_window.py} & Orchestrateur principal \\ \hline
    \texttt{top\_bar.py} & Barre du haut (date, réseau, thème, périodes) \\ \hline
    \texttt{sidebar.py} & Navigation gauche (programmes, classes) \\ \hline
    \texttt{group\_panel.py} & Stats groupe : KPIs, charts, historique \\ \hline
    \texttt{class\_panel.py} & Grille cartes élèves \\ \hline
    \texttt{student\_detail.py} & Détail élève : photo, infos, événements \\ \hline
    \texttt{data\_loader.py} & Requêtes DB (33 méthodes) \\ \hline
    \texttt{event\_actions.py} & CRUD événements + menu contextuel \\ \hline
    \texttt{event\_generator.py} & Wizard génération événement \\ \hline
    \texttt{timetable\_editor.py} & Éditeur emploi du temps \\ \hline
\end{tabular}
\end{center}

\section{Dépendances}

\begin{itemize}
  \item \textbf{PySide6} $\ge$ 6.5 — Framework Qt6
  \item \textbf{psycopg2-binary} $\ge$ 2.9 — Driver PostgreSQL
  \item \textbf{materialyoucolor} $\ge$ 3.0 — Moteur de couleurs Material Design~3
  \item \textbf{LarcCommon} — Librairie partagée (thèmes, i18n, icônes, widgets)
\end{itemize}

\section{Connexion à la base de données}

\begin{itemize}
  \item \textbf{Intranet} : PostgreSQL \texttt{127.0.0.1:5432} — \texttt{NewLarcDB}
  \item \textbf{Cloud} : Supabase \texttt{aws-1-eu-north-1.pooler.supabase.com:6543}
  \item \texttt{autocommit = True} sur toutes les connexions
\end{itemize}

\end{document}
